\section*{Capítulo \ref{ch:g10:layered-ocean} --- El océano en capas}

\begin{solution}{exo:g10:layered-ocean:1}
La \omterm{def:g10:layered-ocean:mixed-layer}{capa de mezcla} es la capa superficial del océano en la que la temperatura
y la salinidad apenas cambian con la profundidad; la \omterm{def:g10:layered-ocean:thermocline}{termoclina} es la capa
situada debajo, en la que la temperatura baja deprisa con la profundidad. El
viento y el oleaje remueven la \omterm{def:g10:layered-ocean:mixed-layer}{capa de mezcla} (con la ayuda, en invierno, del
enfriamiento de la superficie).
\end{solution}

\begin{solution}{exo:g10:layered-ocean:2}
El \omterm{def:g10:layered-ocean:upwelling}{afloramiento} es la subida de agua fría desde debajo de la capa
superficial, que sustituye al agua superficial apartada por el viento. El
agua que sube lleva nutrientes minerales a la capa iluminada, donde el
fitoplancton los usa para crecer; el fitoplancton alimenta al resto de la
cadena alimentaria.
\end{solution}

\begin{solution}{exo:g10:layered-ocean:3}
Pacífico tropical: unos \qty{29}{\degreeCelsius} en la superficie y
\qty{4.5}{\degreeCelsius} a \qty{1000}{m}. Océano Ártico: unos
\qty{-1.4}{\degreeCelsius} en la superficie y cerca de
\qty{0}{\degreeCelsius} a \qty{1000}{m}.
\end{solution}

\begin{solution}{exo:g10:layered-ocean:4}
A \qty{0}{\degreeCelsius} (\qty{1028.1}{kg/m^3} frente a \qty{1023.3}{kg/m^3}
a \qty{25}{\degreeCelsius}). El agua de mar: a \qty{4}{\degreeCelsius}, el
agua pura tiene una densidad de \qty{1000.0}{kg/m^3}, y el agua de mar de
salinidad 35, unos \qty{1027.8}{kg/m^3} (leídos en la gráfica).
\end{solution}

\begin{solution}{exo:g10:layered-ocean:5}
La boya deriva a \qty{1000}{m}, se hunde hasta \qty{2000}{m} y luego sube a
la superficie midiendo por el camino la temperatura, la salinidad (a partir
de la conductividad) y la presión; en la superficie envía su perfil por
satélite y vuelve a hundirse. Las boyas derivan con las corrientes por todo
el océano y funcionan durante años, en todas las estaciones, lejos de las
rutas de los barcos.
\end{solution}

\begin{solution}{exo:g10:layered-ocean:6}
$361.84\times10^6 \times 3.682 \approx \qty{1.33e9}{km^3}$, unos
\num{1.33}~mil millones de kilómetros cúbicos. $100/3682 \approx 0.027$:
menos del \qty{3}{\percent} de la profundidad media.
\end{solution}

\begin{solution}{exo:g10:layered-ocean:7}
$22.7 - 0.5 = \qty{22.2}{\degreeCelsius}$, que se alcanza a \qty{15}{m}: la
\omterm{def:g10:layered-ocean:mixed-layer}{capa de mezcla} tiene unos \qty{15}{m} de profundidad. Entre \qty{20}{m} y
\qty{50}{m}: $(21.8 - 17.3)/30 = \qty{0.15}{\degreeCelsius}$ por metro.
\end{solution}

\begin{solution}{exo:g10:layered-ocean:8}
El sol solo calienta la capa superficial, y el calentamiento y el
enfriamiento estacionales solo se reparten por la \omterm{def:g10:layered-ocean:mixed-layer}{capa de mezcla}: por debajo
de la \omterm{def:g10:layered-ocean:thermocline}{termoclina}, el agua queda aislada de ellos. En invierno la superficie
se enfría, se vuelve más densa y se hunde en el agua de debajo, y las
tempestades agitan el mar: la mezcla llega hondo. En verano el sol calienta
una capa superficial ligera que flota sobre el agua más fría; los vientos,
más suaves, no consiguen empujarla hacia abajo, y la \omterm{def:g10:layered-ocean:mixed-layer}{capa de mezcla} sigue
siendo delgada.
\end{solution}

\begin{solution}{exo:g10:layered-ocean:9}
$0.78/0.18 \approx 4.3$ grados cerca de \qty{10}{\degreeCelsius};
$0.78/0.06 = 13$ grados cerca de \qty{0}{\degreeCelsius}. En los mares
polares fríos, los pequeños cambios de salinidad pesan mucho más en la
densidad que los cambios de temperatura: allí, la salinidad controla la
\omterm{def:g10:layered-ocean:stratification}{estratificación}.
\end{solution}

\begin{solution}{exo:g10:layered-ocean:10}
La temperatura es la misma: el agua profunda es más densa porque es más
salada (33.3 frente a 28.1). La capa que hay entre ellas, donde la salinidad
cambia deprisa con la profundidad, es una \omterm{def:g10:layered-ocean:thermocline}{haloclina}; aquí es también la
\omterm{def:g10:layered-ocean:thermocline}{picnoclina}.
\end{solution}

\begin{solution}{exo:g10:layered-ocean:11}
Empujada hacia abajo, la porción conserva su temperatura y su salinidad, y
por tanto su densidad. Rodeada de agua más densa, es más ligera que el agua
que desplaza: el empuje la devuelve hacia arriba, a su propio nivel. Una
porción empujada hacia arriba sería más densa que su nuevo entorno y volvería
a hundirse. Las capas son estables.
\end{solution}

\begin{solution}{exo:g10:layered-ocean:12}
Masa: $0.001 \times 1028.1 = \qty{1.0281}{kg}$. Nuevo volumen: $1.0281/1023.3
\approx \qty{0.0010047}{m^3}$, unos \qty{1004.7}{mL}: se ha dilatado unos
\qty{4.7}{mL}, un \qty{0.47}{\percent}. La capa de \qty{700}{m} ganaría
$700 \times 0.0047 \approx \qty{3.3}{m}$ de espesor.
\end{solution}

\begin{solution}{exo:g10:layered-ocean:13}
En el Atlántico Norte en febrero. Allí, los \qty{150}{m} superiores tienen
casi la misma temperatura y, por tanto, casi la misma densidad: removerlos
cuesta poco. En el Pacífico tropical, la densidad aumenta $1026.0 - 1021.6 =
\qty{4.4}{kg/m^3}$ entre \qty{50}{m} y \qty{200}{m}: una \omterm{def:g10:layered-ocean:thermocline}{picnoclina} marcada,
y la tempestad tendría que empujar agua ligera hacia abajo, hacia un agua
mucho más densa.
\end{solution}

\begin{solution}{exo:g10:layered-ocean:14}
Mucha clorofila lejos de las costas: la franja a lo largo del ecuador en el
Pacífico (y en el Atlántico), donde los alisios provocan un \omterm{def:g10:layered-ocean:upwelling}{afloramiento}, y
las latitudes altas, como el Atlántico Norte, donde la \omterm{def:g10:layered-ocean:mixed-layer}{capa de mezcla} de
invierno llega hondo y sube nutrientes. Muy poca clorofila: los centros de
los giros subtropicales, muy \omterm{def:g10:layered-ocean:stratification}{estratificados}, donde nada sube el agua
profunda y la capa iluminada se ha quedado sin nutrientes.
\end{solution}

\begin{solution}{exo:g10:layered-ocean:15}
Bajada: $1000/(6 \times 3600) \approx \qty{0.046}{m/s}$, unos
\qty{4.6}{cm/s}. Subida: $2000/(10 \times 3600) \approx \qty{0.056}{m/s}$,
unos \qty{5.6}{cm/s}. $36\,000/2 = 18\,000$ medidas, separadas por término
medio $2000/18\,000 \approx \qty{0.11}{m}$.
\end{solution}

\begin{solution}{pb:g10:layered-ocean:1}
\textbf{1.} El giro, en $24.0 - 19.0 = \qty{5.0}{\degreeCelsius}$.

\textbf{2.} El giro (\qty{20.5}{\degree S} frente a \qty{13.5}{\degree S}).
Al estar más lejos del ecuador, debería ser el más fresco: las latitudes por
sí solas no explican por qué está \qty{5}{\degreeCelsius} más cálido.

\textbf{3.} $24.0 - 0.5 = \qty{23.5}{\degreeCelsius}$, que se alcanza entre
\qty{50}{m} (23.6) y \qty{100}{m} (22.6), hacia $50 + 50 \times
0.1/1.0 = \qty{55}{m}$.

\textbf{4.} No: la temperatura baja de forma regular desde la superficie,
\qty{1.6}{\degreeCelsius} en los primeros \qty{25}{m} y
\qty{5.3}{\degreeCelsius} en los primeros \qty{100}{m}; la \omterm{def:g10:layered-ocean:thermocline}{termoclina} empieza
en la misma superficie.

\textbf{5.} Entre \qty{200}{m} y \qty{300}{m}: $20.3 - 14.4 =
\qty{5.9}{\degreeCelsius}$ por cada cien metros.

\textbf{6.} Frente a Perú, entre \qty{50}{m} y \qty{100}{m}: unos
\qty{70}{m} por interpolación (el perfil completo da \qty{65}{m}). En el
giro, entre \qty{200}{m} y \qty{300}{m}: unos $200 + 100 \times 5.3/5.9
\approx \qty{290}{m}$.

\textbf{7.} Sí, en todos los niveles de las dos columnas: el agua más ligera
descansa sobre el agua más densa, y las dos columnas son estables.

\textbf{8.} El agua más salada debería ser más densa; sin embargo, el agua
superficial del giro es más ligera: \qty{1024.71}{kg/m^3} frente a
\qty{1025.08}{kg/m^3}.

\textbf{9.} El agua del giro es unas 1.3 unidades más salada, lo que añade
unos $1.3 \times 0.8 \approx \qty{1.0}{kg/m^3}$; está
\qty{5}{\degreeCelsius} más cálida, lo que quita unos $5 \times 0.27 \approx \qty{1.4}{kg/m^3}$ a estas temperaturas. Se impone la temperatura.

\textbf{10.} $1026.44 - 1025.08 = \qty{1.36}{kg/m^3}$.

\textbf{11.} $1025.30 - 1024.71 = \qty{0.59}{kg/m^3}$.

\textbf{12.} De la superficie a \qty{500}{m}, la densidad aumenta
$1026.92 - 1024.71 = \qty{2.21}{kg/m^3}$, de los que solo
\qty{0.59}{kg/m^3} están por encima de \qty{200}{m}: la mayor parte,
\qty{1.62}{kg/m^3}, está por debajo de \qty{200}{m}, en la \omterm{def:g10:layered-ocean:thermocline}{termoclina}, entre
\qty{200}{m} y \qty{500}{m}.

\textbf{13.} Frente a Perú, la densidad aumenta $1026.97 - 1025.08 =
\qty{1.89}{kg/m^3}$ de la superficie a \qty{500}{m}, de los que
\qty{1.36}{kg/m^3} están en los \qty{200}{m} superiores: la mayor parte está
cerca de la superficie.

\textbf{14.} A \qty{500}{m}, los dos lugares son casi iguales
(\qty{7.9}{\degreeCelsius} y \qty{7.1}{\degreeCelsius},
\qty{1026.97}{kg/m^3} y \qty{1026.92}{kg/m^3}): la misma agua fría está
debajo de ambos; la diferencia es lo cerca que llega de la superficie.

\textbf{15.} El viento que sopla a lo largo de la costa hacia el ecuador
empuja el agua superficial mar adentro; el agua de debajo sube cerca de la
costa para ocupar su lugar. El agua fría de la \omterm{def:g10:layered-ocean:thermocline}{termoclina} es aspirada hacia
arriba, de modo que la \omterm{def:g10:layered-ocean:thermocline}{termoclina} se inclina hacia la costa y llega casi a la
superficie.

\textbf{16.} Frente a Perú: el agua a entre \qty{13}{\degreeCelsius} y
\qty{15}{\degreeCelsius}, procedente de debajo de la \omterm{def:g10:layered-ocean:thermocline}{termoclina}, está en los
\qty{100}{m} superiores, al alcance de la mezcla del viento y de la luz. En el
giro está por debajo de \qty{250}{m}, bajo una capa cálida y ligera que el
viento no consigue remover hasta tanta profundidad.

\textbf{17.} El \omterm{def:g10:layered-ocean:upwelling}{afloramiento} se detendría: el agua superficial se calentaría
al sol, se formaría encima una capa cálida y ligera, y el aporte de
nutrientes a la capa iluminada iría disminuyendo.

\textbf{18.} La costa de Perú va del azul claro al amarillo (mucha
clorofila); el centro del giro del Pacífico Sur está entre las regiones de
azul más oscuro del mapa (muy poca).

\textbf{19.} El fitoplancton necesita nutrientes además de luz. En el giro,
la fuerte \omterm{def:g10:layered-ocean:stratification}{estratificación} mantiene el agua profunda, rica en nutrientes, muy
por debajo de la capa iluminada, y los nutrientes de la capa superficial se
agotaron hace tiempo.

\textbf{20.} $1.36/0.59 \approx 2.3$: \textbf{el aumento de densidad en los
\qty{200}{m} superiores es unas 2.3 veces mayor frente a Perú que en el
giro}. La \omterm{def:g10:layered-ocean:thermocline}{picnoclina}, con el agua fría y rica en nutrientes que hay debajo,
se ha levantado hasta los \qty{200}{m} superiores frente a Perú, mientras que
en el giro está más honda, por debajo de \qty{200}{m}.
\end{solution}
